
immersion fermée.

Enfin @@M0@@ est un torseur pour f.p.p.f.~; a fortiori, grâce à la descente f.p.q.c. à partir d'un «~splittage~» du torseur, @@M1@@ est représentable par un morphisme affine fidèlement plat de présentation finie (voir~: Oort, «~Commutative group schemes~», Proposition (17.4)).

\textbf{Corollaire 3.3.} \uline{Sous les conditions de} 3.1\uline{, on a}
\begin{enumerate}[label=\roman*)]
\item \uline{Si} @@M2@@ \uline{est un isomorphisme, alors} @@M3@@ \uline{est un isomorphisme.}
\item \uline{En supposant le} @@M4@@\uline{-schéma} @@M5@@ \uline{affine, si} @@M6@@ \uline{est surjectif, alors} @@M7@@ \uline{est surjectif.}
\end{enumerate}

En effet (i) résulte du diagramme suivant de faisceaux pour la topologie de Zariski
@@M8@@
en lui appliquant le lemme des cinq.

(ii) résulte du diagramme à lignes exactes suivant
@@M9@@
car les flèches verticales extérieures sont des isomorphismes, @@M10@@ étant affine.

